Action chunking lets imitation policies execute several actions from one
inference call. The executed chunk length sets a tradeoff: long chunks are
smooth and cheap, short chunks limit how far error compounds before the policy
looks again. Practice picks one chunk length per task. We argue this is the
wrong unit, and study the quantity it trades against directly. For each state
we measure a finite-time maximal Lyapunov exponent by restoring the simulator,
perturbing the first action, replaying the rest open-loop, and fitting the
growth rate of the divergence. The paper then makes three claims. First, a
single task contains both stability regimes: across more than twenty datasets
on three platforms, between 78 and 100 percent of episodes contain both stable
and unstable segments with several crossings each, the unstable segments are
short and concentrated at contact, and the structure is unchanged by a
fourfold change in probe size. Second, the regime can be labeled and predicted:
small heads reading only camera frames and proprioception reach 0.83 to 0.96
AUROC on confident stamps and transfer across tasks, well above a privileged
contact baseline, while a simulator oracle bounds what perfect prediction is
worth. Audited against our labels, the replan signals used by existing
adaptive-horizon methods do not track measured stability, and a kinematic one
is anti-correlated with it. Third, we ask whether the predicted regime can set
chunk length at inference time, and report an incomplete and so far negative
answer: learned predictors gain 3 to 8 points over a fixed 16-step chunk but
stay inside run-to-run noise, never beat a properly swept constant chunk, and a
paired test that forks each episode at a measured unstable moment finds that
closing the loop rescues as many episodes as it breaks. Chunked execution has a
legible and predictable stability structure; converting it into a chunk-length
controller does not yet pay, and we identify the closed-loop axis as the
measurement that would say when it should.
